\documentclass[10pt,a4paper]{amsart}
\usepackage[margin=19mm]{geometry}
\usepackage{amsmath,amssymb,mathtools}
\usepackage[scr=rsfs,cal=euler]{mathalfa}
\usepackage{xfrac}
\usepackage[unicode,hidelinks,bookmarks=false]{hyperref}
\newcommand{\ie}{i.e.\ }
\newcommand{\cf}{cf.\ }
\newcommand{\resp}{respectively\ }
\newcommand{\Subsection}{\subsection}
\providecommand{\nobreakdash}{\nobreak\hbox{-}}
\setlength{\emergencystretch}{2em}
\multlinegap0pt
\usepackage{bm}
\usepackage{bbm}
\usepackage{dsfont}
\usepackage{xspace}
\usepackage{cancel}
\usepackage{mathtools}
\usepackage{amsthm}
\makeatletter
\@ifundefined{theorem}{\newtheorem{theorem}{Theorem}}{}
\makeatother
\title[Invertibility of Anticommutator and Commutators of Higher De]{Invertibility of Anticommutator and Commutators of Higher Degree of $n$-potent Elements}
\author{Vivek Bhabani Lama, Suhas B N}
\date{Source: arXiv:2607.08139v1 (CC BY 4.0)}
\begin{document}
\maketitle

\noindent\emph{Source and licence.} Invertibility of Anticommutator and Commutators of Higher Degree of $n$-potent Elements. Version arXiv:2607.08139v1 (CC BY 4.0), distributed under the Creative Commons Attribution 4.0 International licence, as recorded in the arXiv licence metadata for this exact version and shown on the abstract page https://arxiv.org/abs/2607.08139v1. Full rights notice: licenses/arxiv-2607.08139-CC-BY-4.0.txt.\par\smallskip

\noindent\textit{Editorial scope.} Counted unit: the Theorem whose source label is \texttt{THM: Property Kn for triangular matrix rings} in the article (source0.tex lines 668--677, source label \texttt{THM: Property Kn for triangular matrix rings}), delivered together with the article's own complete original proof of it (source0.tex lines 678--700, one \texttt{proof} environment of 23 lines). No mathematics is written, repaired or re-derived: statement and proof are the article's own text line for line. Required context retained uncounted: none. Every definition, display and label of those lines is retained unchanged. Omitted from this package and not redistributed: 1--667, 701--721. Nothing inside a retained excerpt is omitted or altered: every retained body is delivered exactly as the article's lines are written, so no omission receipt of any kind accompanies this package. Citations: the retained excerpts carry no citation command at all, so no bibliography entry of the article is transcribed and no reference is redistributed. Reference adaptation: none was needed and none was made; every reference inside the delivered unit resolves inside it, so no unresolved reference or citation is produced. Printed slips kept literal: no printed word, symbol, hypothesis or number of the counted statement or of its proof was altered. Layout and class: the article's own class is \texttt{amsart} with packages including amsmath, amssymb, amstext, amsthm, eqnarray, amscd, tikz, tikz-cd, fullpage, float, mathtools, mathabx, graphicx, euscript; the wrapper is the lane's pinned \texttt{amsart} editorial wrapper at 10pt on A4, which restates the theorem-like environments the retained text needs and adds the article's own macro definitions that the retained text needs (none), with extra wrapper packages (bm, bbm, dsfont, xspace, cancel, mathtools). The delivered numbering is the unit's own. Assets: no retained span contains a figure environment, a graphics-inclusion command, a PDF-page inclusion, a table or a TikZ picture; no image file is copied into this package, the package declares no asset dependency and no image file is redistributed with it. Licence: version arXiv:2607.08139v1 of this article is distributed under the Creative Commons Attribution 4.0 International licence, as recorded in the arXiv licence metadata for this exact version and shown on the abstract page https://arxiv.org/abs/2607.08139v1. A full rights notice is delivered at licenses/arxiv-2607.08139-CC-BY-4.0.txt. Characters: every retained excerpt is pure ASCII, so the delivered engine, XeLaTeX with its default Latin Modern text font, reports no missing character.
\begin{theorem}\label{THM: Property Kn for triangular matrix rings}
    Let $S_1$ and $S_2$ be  rings and $M$ be any $(S_1,S_2)$-bimodule and let $T=\begin{pmatrix}
            S_1&M\\
            0&S_2
        \end{pmatrix}$  be the formal triangular ring. Then the following statements hold:
    \begin{enumerate}
        \item For any $n \geq 2$, the ring $T$ satisfies Property $K_n$ if and only if $S_1$ and $S_2$ satisfy Property $K_n$. 
        \item For any $n \geq 2$, the ring $T$ satisfies Property $\overline{K_n}$ if and only if $S_1$ and $S_2$ satisfy Property $\overline{K_n}$. 
    \end{enumerate}
    \end{theorem}

 \begin{proof}
     Let $A=\begin{pmatrix}
         r&m\\
         0&s
     \end{pmatrix} \in T$ be an $n$-potent element. Due to the multiplication operation in $T$, the elements $r$ and $s$ are $n$-potent elements in $S_1$ and $S_2$ respectively. In order to prove the theorem, we give the proof of (1) as (2) follows using similar line of argument. Suppose that $A=\begin{pmatrix}
         r_1 & m_1\\
         0 & s_1
     \end{pmatrix}, B=\begin{pmatrix}
         r_2&m_2\\
         0&s_2
     \end{pmatrix} \in T$ are two $n$-potent elements such that $<A,B>_n$ is a unit in $T$. Using the multiplication in $T$, we have that $<r_1,r_2>_n$ is a unit in $S_1$ and $<s_1,s_2>_n$ is a unit in $S_2$. By our observation, the elements $r_1,r_2$ are $n$-potent in $S_1$ and $s_1,s_2$ are $n$-potent in $S_2$.  Conversely, let $r_1,r_2$ be $n$-potent in $S_1$ and $s_1,s_2$ be $n$-potent in $S_2$ such that $<r_1,r_2>_n$ is a unit in $S_1$ and $<s_1,s_2>$ is a unit in $S_2$. \text{Consider the elements} $ A=\begin{pmatrix}
         r_1 & 0\\
         0 & s_1
     \end{pmatrix} \text{ and } B=\begin{pmatrix}
         r_2&0\\
         0&s_2
     \end{pmatrix} \in T$, clearly both $A$ and $B$ are $n$-potent elements in $T$. Also, we have the following expression for $<A,B>_n$,
     $$<A,B>_n=\begin{pmatrix}
         <r_1,r_2>_n &0\\
         0&<s_1,s_2>_n
     \end{pmatrix} \in T$$
Since $<A,B>_n$ is a unit in $T$, $<r_1,r_2>_n$ is a unit in $S_1$ and $<s_1,s_2>_n$ is a unit in $S_2$. Therefore, $T$ satisfies Property $K_n$ if and only if $S_1$ and $S_2$ individually satisfy Property $K_n$. 
 \end{proof}

\end{document}
